\chapter*{Preface}
\addcontentsline{toc}{chapter}{Preface}

This is an introductory book on linear algebra. It is written for students who
need the subject as a language for mathematics, computation, scientific
modeling, and data. It is also written with future work in artificial
intelligence and machine learning in mind, but it is not a book about neural
networks, transformers, or modern machine-learning architecture.

That distinction is deliberate. Linear algebra prepares students for AI and ML
not by naming every current application, but by teaching the durable habits:
see data as vectors, models as maps, features as coordinates, learning as
approximation, dimension reduction as projection, and low-rank structure as
compression.

The book has two legs. The first is theoretical: definitions, proofs,
invariants, and theorems. The second is computational: algorithms, numerical
experiments, and small programs. Neither leg is optional. A proof without
computation can feel inert; computation without structure can become button
pushing. The point is to learn how the two reinforce each other.

\begin{warningbox}
The phrase ``primed for AI/ML'' means that examples and exercises are chosen to
build useful instincts for data and computation. It does not mean that the
course assumes prior machine-learning knowledge, nor that every topic is
framed through a modern application.
\end{warningbox}
